\documentclass[12pt,a4paper]{article}
\usepackage[english]{babel}
\usepackage{graphicx}
\usepackage{hyperref}
\usepackage{xcolor}

\title{Bachelor Project Writing Guide}
\author{}
\date{}

\begin{document}

\pandocbounded{\includegraphics[keepaspectratio]{media/image11.png}}

Bachelor Project Writing Guide

\begin{quote}
\includegraphics[width=3.60417in,height=0.6175in]{media/image21.png}\emph{\textbf{Panduan
Penulisan Projek Sarjana Muda}}
\end{quote}

\textbf{{[} version 2025 {]}}

\begin{quote}
\textbf{JAWATANKUASA PSM \& PD 2025}

\textsc{MASHANUM OSMAN}

\textsc{AMIRUL RAMZANI RADZID}

\textsc{CHE KU NURAINI CHE KU MOHD}

\textsc{IKMAL FAIQ ALBAKRI MUSTAFA ALBAKRI}

\textsc{MOHD NAZRIEN ZARAINI}

\textsc{MUHAMMAD SHARILAZLAN SALLEH}

\textsc{YAHYA IBRAHIM}

\textsc{ROSMIZA WAHIDA ABDULLAH}
\end{quote}

\ul{TABLE OF CONTENT}

\hyperref[_bookmark0]{PART I : GENERAL GUIDELINES FOR REPORT WRITING 5}

\hyperref[_bookmark1]{{[} ENGLISH VERSION {]} 6}

\hyperref[_bookmark2]{CHAPTER 1. GENERAL GUIDE TO WRITING A REPORT 7}

\hyperref[_bookmark3]{CHAPTER 2. GUIDE TO PREPARING A REPORT 14}

\hyperref[_bookmark4]{CHAPTER 3. GUIDE TO WRITING REFERENCES (HARVARD
STYLE) 18}

\hyperref[_bookmark5]{CHAPTER 4. GUIDE TO WRITING NOTES AND FOOTNOTES
30}

\hyperref[_bookmark6]{{[} BAHASA MELAYU VERSION {]} 31}

\hyperref[_bookmark7]{CHAPTER 1. PANDUAN UMUM UNTUK MENULIS LAPORAN 32}

\hyperref[_bookmark8]{CHAPTER 2. PANDUAN UNTUK MENYEDIAKAN LAPORAN 40}

\hyperref[_bookmark9]{CHAPTER 3. PANDUAN MENULIS RUJUKAN (GAYA
\emph{HARVARD}) 44}

\hyperref[_bookmark10]{CHAPTER 4. PANDUAN PENULISAN NOTA DAN NOTA KAKI
56}

\hyperref[_bookmark11]{PART II : BACHELOR PROJECT CONTENTS GUIDE 57}

\hyperref[_bookmark12]{BITC : Computer Networking 58}

\hyperref[_bookmark13]{Project Type I: System Development 58}

\hyperref[_bookmark14]{CHAPTER 1. INTRODUCTION 59}

\hyperref[_bookmark15]{CHAPTER 2. LITERATURE REVIEW 62}

\hyperref[_bookmark16]{CHAPTER 3. PROJECT METHODOLOGY 64}

\hyperref[_bookmark17]{CHAPTER 4. ANALYSIS AND DESIGN 65}

\hyperref[_bookmark18]{CHAPTER 5. IMPLEMENTATION 68}

\hyperref[_bookmark19]{CHAPTER 6. TESTING 70}

\hyperref[_bookmark20]{CHAPTER 7. PROJECT CONCLUSION 72}

\hyperref[_bookmark21]{REFERENCES 74}

\hyperref[_bookmark22]{BIBLIOGRAPHY 74}

\hyperref[_bookmark23]{APPENDICES 74}

\hyperref[_bookmark24]{PROJECT TYPE II: Analysis (SIMULATION AND
TESTBED) 75}

\hyperref[_bookmark25]{CHAPTER 1. INTRODUCTION 76}

\hyperref[_bookmark26]{CHAPTER 2. LITERATURE REVIEW 79}

\hyperref[_bookmark27]{CHAPTER 3. PROJECT METHODOLOGY 81}

\hyperref[_bookmark28]{CHAPTER 4. DESIGN 83}

\hyperref[_bookmark29]{CHAPTER 5. IMPLEMENTATION 85}

\hyperref[_bookmark30]{CHAPTER 6. TESTING AND ANALYSIS 86}

\hyperref[_bookmark31]{CHAPTER 7. PROJECT CONCLUSION 87}

\hyperref[_bookmark32]{REFERENCES 89}

\hyperref[_bookmark33]{BIBLIOGRAPHY 89}

\hyperref[_bookmark34]{APPENDICES 89}

\hyperref[_bookmark35]{BITD: Database Management 90}

\hyperref[_bookmark36]{CHAPTER 1. INTRODUCTION 91}

\hyperref[_bookmark37]{CHAPTER 2. PROJECT METHODOLOGY AND PLANNING 93}

\hyperref[_bookmark38]{CHAPTER 3. ANALYSIS 95}

\hyperref[_bookmark39]{CHAPTER 4. DESIGN 97}

\hyperref[_bookmark40]{CHAPTER 5. IMPLEMENTATION 99}

\hyperref[_bookmark41]{CHAPTER 6. TESTING 100}

\hyperref[_bookmark42]{CHAPTER 7. CONCLUSION 102}

\hyperref[_bookmark43]{REFERENCES 103}

\hyperref[_bookmark44]{BIBLIOGRAPHY 103}

\hyperref[_bookmark45]{APPENDICES 103}

\hyperref[_bookmark46]{BITE : Game technology 104}

\hyperref[_bookmark47]{CHAPTER 1. INTRODUCTION 105}

\hyperref[_bookmark48]{CHAPTER 2. LITERATURE RIVIEW AND PROJECT
METHODOLOGY 107}

\hyperref[_bookmark49]{CHAPTER 3. ANALYSIS 109}

\hyperref[_bookmark50]{CHAPTER 4. DESIGN 111}

\hyperref[_bookmark51]{CHAPTER 5. IMPLEMENTATION 114}

\hyperref[_bookmark52]{CHAPTER 6. TESTING 117}

\hyperref[_bookmark53]{CHAPTER 7. CONCLUSION 119}

\hyperref[_bookmark54]{REFERENCES 120}

\hyperref[_bookmark55]{BIBLIOGRAPHY 120}

\hyperref[_bookmark56]{APPENDICES 120}

\hyperref[_bookmark57]{BITI : Artificial Intelligence 121}

\hyperref[_bookmark58]{Project Type I: Analysis 121}

\hyperref[_bookmark59]{Artificial Intelligence 122}

\hyperref[_bookmark60]{(BITI) 122}

\hyperref[_bookmark61]{Project Type II: Product 122}

\hyperref[_bookmark62]{CHAPTER 1. INTRODUCTION 123}

\hyperref[_bookmark63]{CHAPTER 2. LITERATURE REVIEW AND PROJECT
METHODOLOGY 125}

\hyperref[_bookmark64]{CHAPTER 3. ANALYSIS 128}

\hyperref[_bookmark65]{CHAPTER 4. DESIGN/THE PROPOSED TECHNIQUE 130}

\hyperref[_bookmark66]{CHAPTER 5. IMPLEMENTATION 134}

\hyperref[_bookmark67]{CHAPTER 6. TESTING 136}

\hyperref[_bookmark68]{CHAPTER 7. CONCLUSION 139}

\hyperref[_bookmark69]{REFERENCES 140}

\hyperref[_bookmark70]{BIBLIOGRAPHY 140}

\hyperref[_bookmark71]{APPENDICES 140}

\hyperref[_bookmark72]{BITM : Media Interactive 141}

\hyperref[_bookmark73]{CHAPTER 1. INTRODUCTION 142}

\hyperref[_bookmark74]{CHAPTER 2. LITERATURE REVIEW AND PROJECT
METHODOLOGY 144}

\hyperref[_bookmark75]{CHAPTER 3. ANALYSIS 147}

\hyperref[_bookmark76]{CHAPTER 4. DESIGN 152}

\hyperref[_bookmark77]{CHAPTER 5. IMPLEMENTATION 155}

\hyperref[_bookmark78]{CHAPTER 6. TESTING 157}

\hyperref[_bookmark79]{CHAPTER 7. CONCLUSION 159}

\hyperref[_bookmark80]{REFERENCES 160}

\hyperref[_bookmark81]{BIBLIOGRAPHY 160}

\hyperref[_bookmark82]{APPENDICES 160}

\hyperref[bits-software-development]{BITS : Software Development 161}

\hyperref[chapter-1.-introduction]{CHAPTER 1. INTRODUCTION 162}

\hyperref[chapter-2.-literature-review-and-project-methodology]{CHAPTER
2. LITERATURE REVIEW AND PROJECT METHODOLOGY 164}

\hyperref[chapter-3.-analysis]{CHAPTER 3. ANALYSIS 166}

\hyperref[chapter-4.-design]{CHAPTER 4. DESIGN 168}

\hyperref[chapter-5.-implementation]{CHAPTER 5. IMPLEMENTATION 171}

\hyperref[chapter-6.-testing]{CHAPTER 6. TESTING 173}

\hyperref[chapter-7.-conclusion]{CHAPTER 7. CONCLUSION 175}

\hyperref[references]{REFERENCES 176}

\hyperref[bibliography]{BIBLIOGRAPHY 176}

\hyperref[appendices]{APPENDICES 176}

\hyperref[part-iii-list-of-appendices]{PART III : LIST OF APPENDICES
177}

\protect\phantomsection\label{_bookmark0}{}

\subsection{\texorpdfstring{\textsc{BITS : Software
Development}}{BITS : Software Development}}\label{bits-software-development}

\subsection{}\label{section}

\subparagraph{CHAPTER 1. INTRODUCTION}\label{chapter-1.-introduction}

Introduction

\begin{quote}
Introduction to your project and your project background as a whole but
in brief.
\end{quote}

Problem statement(s)

\begin{itemize}
\item
  Description of problems that directly influence the motives of the
  project. (Example, explanation of problems of XYZ \& Co. record
  keeping)
\item
  Applicable for project of improving or solving a specific case.
\item
  Project that is of new creation, re-creation, individual initiative or
  not specifically involves any case to deal with, the problems can be
  included in the preamble paragraphs.
\end{itemize}

Objective

\begin{quote}
Can be in point form format (use bullets) together with a brief
explanation.
\end{quote}

Scope

\begin{itemize}
\item
  Describes every scope involved in your project and give reason(s) for
  the involvement. Examples: (e.g. children education), specific users
  (e.g. children between 5 and 8), other specific entities, specific
  platform (e.g. network, OS) etc.
\item
  Can be in point form format (use bullets) together with a brief
  explanation.
\end{itemize}

Project Significance

\begin{quote}
Describes who/what may benefits from the project and how? You can infer
it back to your objectives (if applicable).
\end{quote}

Expected Output

\begin{quote}
Describes what do you expected from your project.
\end{quote}

Conclusion

\begin{quote}
Give a summary of this chapter and next activities to be developed.
\end{quote}

\subparagraph{CHAPTER 2. LITERATURE REVIEW AND PROJECT
METHODOLOGY}\label{chapter-2.-literature-review-and-project-methodology}

Introduction

\begin{quote}
Preview to the literature review and project methodology.
\end{quote}

Facts and findings (based on topic)

\begin{enumerate}
\def\labelenumi{\arabic{enumi}.}
\item
  \textbf{Domain}
\end{enumerate}

\begin{quote}
Identify domain related with your project with explanations.
\end{quote}

Existing System

\begin{itemize}
\item
  Identify domain related with your project with explanations.
\item
  Discuss and state your approach and related or past research,
  references, case study and other finding that relate to your project
  title. Examples: readings 1 and 3, experiments you get 3 and 4, case
  study 5 and 6 etc.
\item
  Tagged the source if you refer the approach from published materials.
\item
  Support the approach with statements (from published materials) to
  justify you fact findings are sound.
\item
  Identify hardware and software used
\end{itemize}

Technique

\begin{itemize}
\item
  State other approaches than what you use that you think also
  applicable and related and give reason to justify why not.
\end{itemize}

Project Methodology

\begin{itemize}
\item
  Describe the selected approach or methodology used in your project
  (Examples: SSADM/SDLC/OOAD, Instructional design, Database life cycle
  and relevant methodology)
\item
  Describe the activities that you may do in every stage and relate with
  your project.
\item
  Include statements (from published materials) that support the
  approach you apply.
\end{itemize}

Project Requirements

\begin{enumerate}
\def\labelenumi{\arabic{enumi}.}
\item
  \textbf{Software Requirement}

  \begin{itemize}
  \item
    List software requirements in point form. Examples: MS Visual Ba-
    sic Professional v.6.x, Apache Tomcat and the related for developing
    application software, MS Project 2000 for project management etc.
  \end{itemize}
\end{enumerate}

Hardware Requirement

\begin{itemize}
\item
  List hardware requirements in point form. Examples: PCs, server,
  devices and storage.
\end{itemize}

Other Requirements

\begin{itemize}
\item
  State other requirements to be used in your project such as you need
  special lab for project development, photo copy facility, discussion
  room and etc.
\end{itemize}

Project Schedule and Milestones

\begin{itemize}
\item
  Explain your actions plan prior to the end of the project. Apply from
  what you have learnt from project management.
\item
  List and describe stage by stage of your activities.
\item
  Attach your project time line or Gantt chart (in 1 page view)
\end{itemize}

Conclusion

\begin{quote}
Summarize the chapter and explain the next activities to be developed.
\end{quote}

\subparagraph{CHAPTER 3. ANALYSIS}\label{chapter-3.-analysis}

Introduction

\begin{quote}
Preview to the analysis phase and how it would be developed.
\end{quote}

Problem Analysis

\begin{itemize}
\item
  Investigate and describe current system scenario/situation. Reiterate
  the data flow diagram or activity diagram from your reference(s),
  showing how the current system(s) or business(s) runs.
\item
  Use appropriate diagram to visualize the system flow such as sequence
  diagram (if use OOAD with UML), DFD (if use SSADM) and Flowchart (for
  network project).
\item
  For simulation based project, analyse the real graph given by company.
\item
  Explain in detail the problems statement as mentioned in Chapter 1.
\end{itemize}

Requirement analysis

\begin{enumerate}
\def\labelenumi{\arabic{enumi}.}
\item
  \textbf{Data Requirement}

  \begin{itemize}
  \item
    What data should the system input and output, and what data should
    the system store internally.
  \item
    It can be illustrated by using Data Model or Data Dictionary.
  \end{itemize}
\end{enumerate}

Functional Requirement

\begin{itemize}
\item
  Specify the functions of the system, how it records, compute, trans-
  forms, and transmits data.
\item
  It can be illustrated by using DFD, Context diagram or Use case.
\end{itemize}

Non-functional Requirement

\begin{quote}
Specify how well the system performs its intended functions. E.g.
Quality requirement, Performance-how many computer resources should it
use, how accurate should the result be, how much data should it be able
to store?
\end{quote}

Others Requirement

\begin{quote}
Describe each of software, hardware and requirements that will be used
(justification of usage).
\end{quote}

Conclusion

\begin{quote}
Summarize the chapter and explain the next activities to be developed.
\end{quote}

\subparagraph{CHAPTER 4. DESIGN}\label{chapter-4.-design}

Introduction

\begin{quote}
Introductory preview to this chapter. Examples: this chapter defines the
results of the analysis of the preliminary design and the result of the
detailed design.
\end{quote}

High-Level Design

\begin{quote}
Describe high-level view of your system's structure or system's
interior. You may refine some of the function below once you are in PSM
II.
\end{quote}

System Architecture

\begin{itemize}
\item
  Define architecture view of your system. The architecture view can be
  presented in layer, framework, tier or patent (for OOAD).
\item
  Describe the static view and dynamic view of your application (for
  OOAD - use interaction diagram, high-level class diagram).
\item
  For those use SSADM: use any appropriate diagram to view your system.
\end{itemize}

User Interface Design

\begin{quote}
Refine your user interface design that defined in chapter 5 of PSM I.
\end{quote}

\begin{enumerate}
\def\labelenumi{\alph{enumi}.}
\item
  Navigation Design
\end{enumerate}

\begin{quote}
Define and refine the navigation flow and types of navigation controls.
\end{quote}

\begin{enumerate}
\def\labelenumi{\alph{enumi}.}
\setcounter{enumi}{1}
\item
  Input Design

  \begin{itemize}
  \item
    Define and refine the screens (e.g. types of inputs such as text,
    numbers, and selection box etc.) used to enter information, as well
    as any forms on which users write or type information.
  \item
    Define and refine validation rule for each of input field.
  \end{itemize}
\item
  Output Design
\end{enumerate}

\begin{quote}
Define and refine the types of outputs including detail reports, summary
reports, turnaround documents and graphs. Classify your output in term
of periodically or ad-hoc basis. For example daily, monthly, yearly etc.
\end{quote}

Database Design

\begin{enumerate}
\def\labelenumi{\arabic{enumi}.}
\item
  \textbf{Conceptual and Logical Database Design}

  \begin{itemize}
  \item
    Introduce briefly to the logical data model (LDM) or entity
    relation- ship diagram (ERD) and what they have to do with your
    database design.
  \item
    Define, refine and construct the entity relationship (ER) diagrams
    in details with explanation in text on what basis (business rules)
    you apply for every entity relationship you have defined in the
    diagram.
  \item
    Data dictionary and normalization
  \end{itemize}
\end{enumerate}

Detailed Design

\begin{quote}
Specification and diagrams may be further elaborated. Emphasis should be
on the logic of the design and the approach to satisfying the
requirements -- How will the system function?
\end{quote}

Software Design

\begin{itemize}
\item
  For those use SSADM/SDLC: Describe in detail of every functions
  according to your DFD in format of program specification. The program
  specification includes information of program description, file
  input/output, pseudo code and attach sample screens.
\item
  For those use OOAD/UML: Describe in detail of every classes. Examples:
  class name, responsibility, attributes methods/operations. Describe
  every methods/operations such as its responsibility, input/output
  parameter, pre/post condition and algorithm --should be bias to
  language chosen for your software development (e.g. in English like of
  Java code).
\end{itemize}

Physical Database Design

\begin{quote}

\end{quote}

\end{document}
Translate logical to target DBMS (use DDL/DCL) - base tables, de- sign
for other business rules (constraint- validation for fields, records),
file organization and indexes
\end{quote}

Conclusion

\begin{quote}
Summarize the chapter and explain the next activities to be developed
\end{quote}

\subparagraph{CHAPTER 5.
IMPLEMENTATION}\label{chapter-5.-implementation}

Introduction

\begin{itemize}
\item
  Introductory preview to this chapter.
\item
  Briefly describe the activity involved in the implementation phase and
  what is the expected output after you complete this phase. Provide
  chapter outline diagram of Chapter 5.
\end{itemize}

Software Development Environment setup

\begin{itemize}
\item
  Define your development environment setup.
\item
  Use diagram to view/present the environment architecture. Examples:
  deployment diagram, draw the software (client s/w, server s/w),
  hardware (server configuration e.g. port no., IP address, database
  instance name, table space etc.) and network setup.
\end{itemize}

Software Configuration Management

\begin{enumerate}
\def\labelenumi{\arabic{enumi}.}
\item
  \textbf{Configuration environment setup}

  \begin{itemize}
  \item
    Explain how you design and setup the configuration management in
    your project. Do not explain on how to install the software.
  \item
    Explain (if any) the software tools used to support your
    configuration control. E.g. Visual SourceSafe, CVS etc.
  \end{itemize}
\end{enumerate}

Version Control Procedure

\begin{quote}
Describe the procedure and control in managing your source code version.
\end{quote}

Implementation Status

\begin{quote}
Describe the progress of the development status for each of the
component/module. For example: component/module name, description,
duration to complete, date completed, size of software etc.
\end{quote}

Conclusion

\begin{quote}
Summarize the chapter and explain the next activities to be developed.
\end{quote}

\subparagraph{CHAPTER 6. TESTING}\label{chapter-6.-testing}

Introduction

\begin{quote}
Introductory preview to this chapter. E.g. briefly describe the activity
involved in testing phase and what is the testing strategy to be adopted
in your project. Provide chapter outline diagram of Chapter 6.
\end{quote}

Test Plan

\begin{enumerate}
\def\labelenumi{\arabic{enumi}.}
\item
  \textbf{Test Organization}
\end{enumerate}

\begin{quote}
Describe personnel involved.
\end{quote}

Test Environment

\begin{itemize}
\item
  Describe the location/environment of testing to be carried out.
\item
  Define hardware, firmware configurations, preparations and training
  prior to testing.
\end{itemize}

Test Schedule

\begin{quote}
Define how many cycles and duration of your test to be conducted.
\end{quote}

Test Strategy

\begin{quote}
Explain the strategy to be selected such as bottom-up or top-down and
black-box/white-box classes of tests.
\end{quote}

Classes of tests

\begin{quote}
Output correctness or functionality test, Security test, Stress test and
etc.
\end{quote}

Test Design

\begin{enumerate}
\def\labelenumi{\arabic{enumi}.}
\item
  \textbf{Test Description}
\end{enumerate}

\begin{quote}
Test case identification, test cases and expected result for each module
are designed and documented.
\end{quote}

Test Data

\begin{quote}
Real life or synthetic data will be selected.
\end{quote}

Test Results and Analysis

\begin{itemize}
\item
  Test case identification, tester identification, test case results
  (Success/Fail), and detailed documentation on the failed test case.
\item
  How satisfied overall your intended users and yourself with the
  system.
\end{itemize}

Conclusion

\begin{quote}
Summarize the chapter and explain the next activities to be developed.
\end{quote}

\subparagraph{CHAPTER 7. CONCLUSION}\label{chapter-7.-conclusion}

Observation on Weaknesses and Strengths

\begin{itemize}
\item
  State the weaknesses and strength of your project.
\item
  You also may state other's responses regarding project topics.
\end{itemize}

Propositions for Improvement

\begin{itemize}
\item
  Present your suggestions on how your system can be improved better.
\item
  Elaborate each of your suggestions in paragraph.
\end{itemize}

Project Contribution

\begin{itemize}
\item
  State your project contribution to the
  university/faculty/company/individual.
\item
  State where to find the user manual - e.g. Appendix XX
\end{itemize}

Conclusion

\begin{itemize}
\item
  State whether you think your project meets your set objectives
  conclusively.
\item
  Concluding phrases to conclude the project.
\end{itemize}

\subparagraph{REFERENCES}\label{references}

\begin{itemize}
\item
  A list of references used in the project report such as journals,
  articles, books and etc. as directly quoted in the report.
\item
  For complete reference, please refer to Final Year Project Writing
  Guide 2025.
\end{itemize}

\subparagraph{BIBLIOGRAPHY}\label{bibliography}

\begin{itemize}
\item
  A list of references used in the project report such as journals,
  articles, books and etc. but not directly quoted in the report.
\item
  For complete reference, please refer to Final Year Project Writing
  Guide 2025.
\end{itemize}

\subparagraph{APPENDICES}\label{appendices}

\begin{itemize}
\item
  E.g. user guide, user manual, diagrams and etc.
\item
  The Turnitin report for the first page should be included in the
  appendix.
\end{itemize}

\section{PART III : LIST OF
APPENDICES}\label{part-iii-list-of-appendices}

\begin{quote}
\textbf{IMPORTANT NOTES:}

The following appendices are purposely for clearer understanding on how
does your report look like.

\textbf{Please kindly to download and use the provided template version
from the course page in the e-Learning platform.}
\end{quote}
